We survey wave-type models from engineering in which finite-time singularities
form with self-similar structure. The models are electrostatic MEMS with
inertia, thermal explosion under Maxwell--Cattaneo heat conduction, strings
and membranes on softening foundations ($\varphi^4$/Klein--Gordon), nonlinear
acoustics, director fields (wave maps) and relativistic membranes. For each
we analyse the stability of the self-similar blowup with the method of
Donninger and coauthors (similarity coordinates, semigroup theory, spectral
analysis, modulation). For every ODE-type blowup with scaling exponent $s>-1$,
the linearisation turns out to have a universal spectrum
$\{1-\ell-2n\}\cup\{-2s-2-\ell-2n\}$, whose only non-negative points come from
time translation and Lorentz boosts. For the hyperbolic MEMS equation
$w_{tt}+\gamma w_t-\Delta w=-\lambda w^{-q}$, $q>1$, in dimensions $d\le3$, this yields (with full
spectral analysis and a proof sketch of the nonlinear step) that touchdown
$w\simeq c(T-t)^{2/(q+1)}$ with a flat profile in the light cone is stable
modulo symmetries, also in the presence of damping, with rate
$(T-t)^{\omega}$, $\omega<\min\{1,2(q-1)/(q+1)\}$. Non-trivial focused touchdown
profiles are found numerically in dimensions $2$ and $3$; numerically they
carry additional unstable modes. We
also identify where the method fails or the model changes the singularity
type: the full Cattaneo model has a parabolic blowup rate, Arrhenius kinetics
prevents blowup, and quasilinear acoustic and membrane models remain open.
The spectra at the ODE profiles are cross-checked by high-precision collocation
and, in selected cases, Evans-function tests. Direct simulations of the clamped MEMS device confirm
the touchdown exponent, the constant $c$ and the convergence rate $2/3$. They
also show that exponents fitted at moderate gaps are biased low.
